The effective interaction, depending from energy scale and momentum, can be attractive or repulsive. In the case of an attractive interaction, the system can be unstable and undergo a phase transition to a new state of matter. The nature of this new state depends on the specific form of the effective interaction and the symmetries of the system. In other cases, it can lead to charge density waves, spin density waves, or other types of ordered states. The study of effective interactions is crucial for understanding the rich variety of phases and phenomena observed in condensed matter systems.

The divergence is not very physical and it can be removed by considering the finite size of the system or by introducing a cutoff in the momentum space.

\[
    V_{eff}(q,i \nu_n) = Draw
\]

[\textit{Draw of sum of wiggle and spring to obtain zigzag to represent the effective interaction}]

Now we can consider the rpa approximation, which is a way to sum up an infinite series of diagrams in perturbation theory. In the rpa approximation, we consider the bubble diagrams, which represent the polarization of the system due to the interactions. The effective interaction in the rpa approximation can be written as:
\[
    V^{eff}_{RPA}(q,i \nu_n) = [\textit{zigzag + zigazag-bubble-zigzag + zigzag-bubble-zigzag-bubble-zigzag + ... = boldzigzag}]
\]
remembering
\[
    V^{eff}(q,i \nu_n) = V_c(q) \frac{(i \nu_n)^2}{(i \nu_n)^2 - \Omega^2}
\]
we call all that series a bold zigzag, and we can write it as:
\[
    \left(1 + [\textit{zigzag-bubble}]\right)V_{RPA}^{eff}(q,i \nu_n) = [\textit{zigzag}] = V^{eff}(q,i \nu_n),
\]
since it is a sort of geometric series, we can write it as:
\[
    V^{eff}_{RPA}(q,i \nu_n) = \frac{V^{eff}(q,i \nu_n)}{1 - V^{eff}(q,i \nu_n) \Pi_0(q,i \nu_n)} = \frac{V_c(q)}{1-V_c(q) \Pi_0(q,i \nu_n)}\frac{(i \nu_n)^2}{(i \nu_n)^2 - \frac{\Omega^2}{1-V_c(q) \Pi_0(q,i \nu_n)}},
\]
where we can recognize the renormalized phonon frequency:
\[
    \frac{\Omega}{\sqrt{1 - V_c(q) \Pi_0(q,i \nu_n)}},
\]
and the screened Coulomb interaction:
\[
    \frac{V_c(q)}{1 - V_c(q)\Pi_0(q,i \nu_n)},
\]
and now we can rewrite the effective interaction in the rpa approximation as:
\[
    \dots
\]
and since the renormalized phonon frequency has an imaginary part, it can be used to remove the divergence of the effective interaction, and it can lead to a finite value of the effective interaction at low energy and momentum.

    [\textit{Draw of effective potential with rpa approximation, in which the curve do not goes to infinity and omegaq}]

The main message is that the effective interaction between two electrons in a lattice, mediated by phonons, can be attractive at low energy and momentum, and it can lead to the formation of Cooper pairs and superconductivity. The rpa approximation is a powerful tool to study the effective interaction in condensed matter systems, and it can provide insights into the mechanisms of superconductivity and other phenomena.

The new frequency
\[
    \omega_q = \frac{\Omega}{\sqrt{1 - \frac{4\pi}{q^2} \Pi_0(q,i \nu_n)}} \simeq \frac{\Omega}{\sqrt{1 - \frac{4\pi}{q^2} \mathcal{N}(0)}} \xrightarrow{q \to 0} \frac{\Omega}{\sqrt{1 + \frac{q_{TF}^2}{q^2}}} \sim \Omega \frac{q}{q_{TF}},
\]
where we have introduced the Thomas-Fermi wavevector:
\[
    q_{TF}^2 = 4\pi \mathcal{N}(0),
\]
and we can see that the phonon frequency goes to zero as $q \to 0$. This means also that the frequancy goes linearly with $q$ at small $q$.


\section{Cooper pairs and superconductivity}

 [...]

 [\textit{Draw of sphere of metal, we add two extra electrons with opposite momenta and spin on top of it, forming a narrow shell (dotted circle) with energy proportional to frequency of }]

fermi energy \(\sim \) 10eV, phonon frequency \(\sim\) 0.01eV, so the shell is very narrow, and the electrons in this shell can form Cooper pairs due to the attractive interaction mediated by phonons. The Cooper pairs can condense into a superconducting state at low temperatures, leading to zero resistance and other remarkable properties of superconductors. The study of Cooper pairs and superconductivity is a fundamental topic in condensed matter physics, and it has important implications for both fundamental science and technological applications.

In order to form a cooper pair, we do not need only a finite attractive interaction, but we also need a finite density of states at the Fermi level, in order to have a finite number of electrons that can form Cooper pairs. There is a fermi sphere, and we add two extra electrons with opposite momenta and spin on top of it, forming a narrow shell with energy proportional to the frequency of the phonons. Then we can consider the attractive potential mediated by phonons.

From a first quantized perspective we want to solve the Schrodinger equation for two electrons interacting via the effective interaction mediated by phonons. The Schrodinger equation can be written as:
\[
    \left[\frac{\hat{P}_1^2}{2m} + \frac{\hat{P}_2^2}{2m} + U(\hat{r}_1 , \hat{r}_2) \right]\psi(r_1, \sigma_1 ; r_2, \sigma_2) = E \psi(r_1, \sigma_1 ; r_2, \sigma_2),
\]
we wnat to find the bound state of two electrons, and get the energy of the bound state.

We can factorize the wavefunction as:
\[
    \psi(r_1, \sigma_1 ; r_2, \sigma_2) = \phi(r_1, r_2) \chi(\sigma_1, \sigma_2),
\]
as usual, in a radial and angular part. The spinorial component can be either a singlet or a triplet, and the spatial component can be either symmetric or antisymmetric. The spinor is given by
\[
    \chi(\sigma_1, \sigma_2) = \begin{dcases}
        (s=0) \quad \frac{1}{\sqrt{2}} \left( \ket{\uparrow \downarrow} - |\downarrow \uparrow \rangle \right) & \text{singlet} \\
        (s=1) \quad \begin{dcases}
                        \ket{\uparrow \uparrow}                                                                    \\
                        \frac{1}{\sqrt{2}} \left( \ket{\uparrow \downarrow} + |\downarrow \uparrow \rangle \right) \\
                        \ket{\downarrow \downarrow}
                    \end{dcases}
    \end{dcases}
\]
The spinorial part can be either symmetric or antisymmetric, and the spatial part can be either symmetric or antisymmetric. The total wavefunction must be antisymmetric under the exchange of the two electrons, so if the spinorial part is a singlet (antisymmetric), the spatial part must be symmetric, and vice versa.

The radial part of the wavefunction can be written as:
\[
    \phi(r_1, r_2) = \sum_{\mathbf{k}} g(\mathbf{k}) \frac{e^{i \mathbf{k} \cdot \mathbf{r}_1 }}{\sqrt{V}} \frac{e^{-i \mathbf{k} \cdot \mathbf{r}_2 }}{\sqrt{V}} = \sum_{\mathbf{k}} g(\mathbf{k}) \frac{e^{i \mathbf{k} \cdot (\mathbf{r}_1 - \mathbf{r}_2) }}{V},
\]
and in order to infer the symmetry of the spatial part, we have to exploit the nature of the summation constant \(g(\mathbf{k})\). If \(g(\mathbf{k})\) is an even function of \(\mathbf{k}\), then the spatial part is symmetric, and if \(g(\mathbf{k})\) is an odd function of \(\mathbf{k}\), then the spatial part is antisymmetric.

Then the total wavefunction can be written as:
\[
    \psi(r_1, \sigma_1 ; r_2, \sigma_2) = \sum_{\mathbf{k}} g(\mathbf{k}) \frac{e^{i \mathbf{k} \cdot (\mathbf{r}_1 - \mathbf{r}_2) }}{V} \chi^S,
\]
where \(\chi^S\) is the spinorial part, which can be either a singlet (\(S=0\)) or a triplet (\(S=1\)). If we exchange the two particles, we get a minus sign on the exponential (we exchanged \(\mathbf{r}_1\) and \(\mathbf{r}_2\)), but we need to add a condition to the function \(g(\mathbf{k})\) in order to have the total wavefunction antisymmetric:
\[
    \begin{aligned}
        S=0) \quad g(\mathbf{k}) & = g(-\mathbf{k}),  \\
        S=1) \quad g(\mathbf{k}) & = -g(-\mathbf{k}),
    \end{aligned}
\]
so if the spinorial part is a singlet, the spatial part must be symmetric, and if the spinorial part is a triplet, the spatial part must be antisymmetric.

The symmetry of the spatial part will become pretty relevant in a moment, when we discuss...

We put the two electrons in a narrow shell around the Fermi surface, such that from this follows a constraint:
\[
    g(\mathbf{k}) = 0 \quad \text{if} \quad E_{\mathbf{k}} = \frac{\hbar^2 k^2}{2m} < E_F , \\
    g(\mathbf{k}) = 0 \quad \text{if} \quad E_{\mathbf{k}} = \frac{\hbar^2 k^2}{2m} > E_F + \hbar \omega_D,
\]
where we know the Debye frequency \(\omega_D \ll E_F\) to be orders of magnitude smaller than the Fermi energy, so the shell is very narrow.

Let us now define
\[
    U_{k k^{\prime}} = \frac{1}{V} \int \mathrm{d}^3 r_1 \mathrm{d}^3 r_2 \, e^{i (\mathbf{k} - \mathbf{k}^{\prime}) \cdot (\mathbf{r}_1 - \mathbf{r}_2)} U(\mathbf{r}_1, \mathbf{r}_2),
\]
and we can write the Schrodinger equation as:
\[
    \left[ - \frac{\hbar^2}{2m} (\nabla_1^2 + \nabla_2^2) + U(\mathbf{r}_1, \mathbf{r}_2) \right] \dots = 0,
\]
such that
\[
    \sum_{k^{\prime}} \left[ 2 \frac{\hbar^2 \vert \mathbf{k}^{\prime} \vert^2 }{2m} + U(\mathbf{r}_1, \mathbf{r}_2) - E\right] g(\mathbf{k}^{\prime}) \frac{1}{V} e^{i(\mathbf{k}^{\prime}- \mathbf{k}) \cdot (\mathbf{r}_1 - \mathbf{r}_2)} \chi^S = 0,
\]
and if we now integrate over \(\mathbf{r}_1\) and \(\mathbf{r}_2\), we get:
\[
    \int \int \mathrm{d}^3 r_1 \mathrm{d}^3 r_2 \, \sum_{k^{\prime}} \left[ 2 \frac{\hbar^2 \vert \mathbf{k}^{\prime} \vert^2 }{2m} + U(\mathbf{r}_1, \mathbf{r}_2) - E\right] g(\mathbf{k}^{\prime}) \frac{1}{V} e^{i(\mathbf{k}^{\prime}- \mathbf{k}) \cdot (\mathbf{r}_1 - \mathbf{r}_2)} \chi^S = 0,
\]
and it will still be zero if the term in square brackets is zero, so we can write:
\[
    \sum_{k^{\prime}} \left[ 2 E_{k^{\prime}} - E \right] g(k^{\prime}) \int \int \mathrm{d}^3 r_1 \mathrm{d}^3 r_2 \, \frac{1}{V} e^{i(\mathbf{k}^{\prime}- \mathbf{k}) \cdot (\mathbf{r}_1 - \mathbf{r}_2)} + \sum_{k^{\prime}} g(k^{\prime}) \int \int \mathrm{d}^3 r_1 \mathrm{d}^3 r_2 \, U(\mathbf{r}_1, \mathbf{r}_2) \frac{1}{V} e^{i(\mathbf{k}^{\prime}- \mathbf{k}) \cdot (\mathbf{r}_1 - \mathbf{r}_2)} = 0,
\]
which has to be true for any \(\chi^S\), and we can recognize the Fourier transform of the potential \(U_{k k^{\prime}}\), while the first term gives basically a delta function, so we can write:
\[
    (2 E_k - E)g(\mathbf{k}) + \sum_{k^{\prime}} U_{k k^{\prime}} g(\mathbf{k}^{\prime}) = 0,
\]
which is an equation for the coefficients \(g(\mathbf{k})\), which we know to be nonzero only in a narrow shell around the Fermi surface:
\[
    E_F < E_k,\,E_{k^{\prime}} < E_F + \hbar \omega_D.
\]
In this narrow shell, we can approximate the potential \(U_{k k^{\prime}}\) as a constant, since it is mediated by phonons, and it is attractive, so we can write:
\[
    U_{k k^{\prime}} = - \vert U_0 \vert,
\]
which is just an hypothesis, since in real materials the potential can have a more complicated form, but it is a good approximation to understand the physics of Cooper pairs and superconductivity. When this interaction is constant and attractive in the narrow shell around the Fermi surface, we can write the equation for \(g(\mathbf{k})\) as:
\[
    (2E_k -E)g(\mathbf{k}) - \vert U_0 \vert \sum_{k^{\prime}} g(\mathbf{k}^{\prime}) = 0.
\]
Now we have to make use of the symmetry of the spatial part of the wavefunction, which is symmetric for the singlet case and antisymmetric for the triplet case. In the triplet case, we can write:
\[
    E = 2E_k,
\]
since the summation goes to zero due to the antisymmetry of the spatial part, and we have just a sum of their energy: it is not a bound state, since the energy is just the sum of the two electrons, and it is not lower than the energy of two free electrons. In the singlet case, we can write:
\[
    E = 2E_k - \vert U_0 \vert N,
\]
where \(K\) is just a constant to highlight that the energy is lower than the energy of two free electrons, and it is a bound state (since the summation of symmetric functions gives a nonzero finite result).

We can write the equation for \(g(\mathbf{k})\) as:
\[
    \sum_{k} g(\mathbf{k}) = \sum_{k} \frac{\vert U_0 \vert \sum_{k^{\prime} }g(\mathbf{k}^{\prime} ) }{2E_k - E},
\]
but since the two summations are equal to \(N\) we can write the expression for the energy which solves our schrodinger equation:
\[
    1 = \vert U_0 \vert \sum_{k} \frac{1}{2E_k - E},
\]
and it is often reffered to as the \textbf{gap equation}, since it gives the energy of the bound state, which is the energy gap of the superconductor. The solution of this equation gives us the energy of the Cooper pair, which is lower than the energy of two free electrons, and it is a bound state. The existence of this bound state is a key ingredient for the formation of Cooper pairs and superconductivity.

We are summing over \(\mathbf{k}\) functions of the energy, so we can use the density of states to write the summation as an integral over energy:
\[
    \sum_{k} \mathcal{F}(E_k) = \sum_{k} \int \mathrm{d}\omega \, \delta(\omega - E_k) \mathcal{F}(\omega) = \int \mathrm{d}\omega \, \left[ \sum_{k} \delta(\omega - E_k) \right] \mathcal{F}(\omega) = \int \mathrm{d}\omega \, D(\omega) \mathcal{F}(\omega),
\]
where the formal expression in square brackets is the definition of the density of states \(D(\omega)\). Then we can write the gap equation as:
\[
    1 = \vert U_0 \vert \int_{E_F}^{E_F + \hbar \omega_D} \mathrm{d}\omega \, \frac{D(\omega)}{2\omega - E},
\]
where we restricted the integration to the narrow shell around the Fermi surface(since even in the starting summation over \(\mathbf{k}\) we were selecting only the momenta which correspond to energies in that shell).

    [\textit{Draw of density of states as a function of energy, with a narrow shell around the Fermi level}]

If we plot for example the density of states of the standard jellium model in 3D, we can see that it is a smooth function of energy, and it is finite at the Fermi level; since the Debye frequency is much smaller than the Fermi energy, we can approximate the density of states as a constant in the narrow shell around the Fermi surface, and we can write:
\[
    1 = \vert U_0 \vert D(E_F) \int_{E_F}^{E_F + \hbar \omega_D} \mathrm{d}\omega \, \frac{1}{2\omega - E},
\]
which can be easily computed knowing that the primitive is a logarithm, and we get:
\[
    1 = \frac{1}{2} \vert U_0 \vert D(E_F) \ln \left( \frac{2E_F + 2\hbar \omega_D - E}{2E_F - E} \right).
\]
Now inverting this equation, we can write the energy of the bound state as:
\[
    \frac{2 E_F - E}{2E_F + 2\hbar \omega_D - E} = e^{-\frac{2}{\vert U_0 \vert D(E_F)}},
\]
and it is now clear why we need to be on top of the fermi surface: if we were in the vacuum we couldn't have made the approximation of a constant density of states, since the density of states in the vacuum is zero at zero energy, and it increases as a power law with energy. In the presence of a finite density of states at the Fermi level, we can obtain a finite value for the energy of the bound state, which is lower than the energy of two free electrons, and it is a key ingredient for the formation of Cooper pairs and superconductivity.

The numerator of the left hand side is the bounding energy of the Cooper pair, which is the energy gap between the energy of the bound state and the energy of two free electrons, and we call it \(\Delta_b\). Thus we can write
\[
    \frac{\Delta_b}{\Delta_b + 2 \hbar \omega_D} = e^{-\frac{2}{\vert U_0 \vert D(E_F)}},
\]
and we can write the energy of the bound state as:
\[
    \Delta_b = \frac{\hbar \omega_D e^{-\frac{1}{\vert U_0 \vert D(E_F)}}}{\sinh(1 / \vert U_0 \vert D(E_F))},
\]
and since \(\sinh(x) \sim x\) for \(x \to 0\), we can write the energy of the bound state as:
\[
    \Delta_b \simeq 2 \hbar \omega_D e^{-\frac{1}{\vert U_0 \vert D(E_F)}},
\]
which is the famous BCS formula for the energy gap of the superconductor, and it shows that the energy gap is exponentially small in the inverse of the product of the attractive interaction and the density of states at the Fermi level. This means that even a weak attractive interaction can lead to a finite energy gap, and it is a key ingredient for the formation of Cooper pairs and superconductivity.

If you in the system there are thermal fluctuations, they can break the Cooper pairs, and they can lead to a finite temperature at which the superconducting state is destroyed. Ideally, starting from an high temperature, we can cool down the system, and at a certain critical temperature \(T_c\) the Cooper pairs will start to form, and the system will undergo a phase transition to a superconducting state. The critical temperature can be estimated from the energy gap, and it is of the order of few Kelvins, since we have an exponential dependence on the inverse of the product of the attractive interaction and the density of states at the Fermi level, which are both small quantities. The thermal energy has to be smaller than the energy gap in order to have a stable superconducting state.

To form a good superconductor, we need a large density of states at the Fermi level, and a strong attractive interaction mediated by phonons. In real materials, the situation is more complicated, since there are many other factors that can affect the superconducting properties, such as impurities, disorder, and other types of interactions. However, if we have a term \(\vert U_0 \vert\) which is even only infinitesimally different from zero, we can have a finite energy gap, and thus a superconducting state. This is the key message of the BCS theory of superconductivity, which is one of the most successful theories in condensed matter physics, and it has been confirmed by many experiments.

We cannot get that term taylor expanding the effective interaction mediated by phonons, since it is zero at zero energy and momentum, but we can get it by considering infinite order perturbation theory, but as presented here it is an analytic result, since we have just made some approximations to solve the Schrodinger equation for two electrons interacting via an attractive potential mediated by phonons. The key ingredient is the existence of a finite density of states at the Fermi level, which allows us to obtain a finite energy gap even for a weak attractive interaction.

Electron-phonon coupling is crucial, because the larger it is the larger the attractive interaction mediated by phonons, and thus the larger the energy gap and the critical temperature of the superconductor. Good conductors (not super-) have a very very small electron-phonon coupling, and thus they do not have a superconducting state (if they had high electron-phonon coupling, they would not conduct well, since the electrons would be scattered by phonons). On the other hand, good superconductors have a large electron-phonon coupling, which leads to a strong attractive interaction mediated by phonons, and thus to a large energy gap and a high critical temperature, but they are not good conductors, since the electrons are strongly scattered by phonons. This is a trade-off between conductivity and superconductivity, and it is one of the reasons why we do not have room temperature superconductors yet.